% \section{Current Consumption Analysis and Battery Selection}
% A practical current consumption analysis has been made to determine a suitable battery capacity and performance.

% \subsection{Theoretical Current Consumption Analysis}
% An estimated current consumption analysis was made based on values from the datasheet of EFR32MG22E \cite{EFR32MG22E}.
% With an approximate number of hours in 10 years being
% \begin{equation}
%     H_{10y} = 24~\text{h/day}\cdot 365~\text{days/year}\cdot 10~\text{years} = 87600~\text{h},
% \end{equation}
% and with wake-ups planned for once a month (except the first month with weekly wake-ups), the estimated total minimum number of wake-ups is
% \begin{equation}
%     N_{WU}=123.
% \end{equation} 

% The following assumptions are used:

% \begin{itemize}
%   \item Wake-up time from EM4: $t_{\text{WU}} = \SI{1.83}{\milli\second}$.
%   \item ADC sample time: $t_{\text{ADC}} = \SI{10}{\milli\second}$.
%   \item TX time: $t_{\text{TX}} = \SI{100}{\milli\second}$ at $I_{\text{TX}} = \SI{8.2}{\milli\ampere}$.
%   \item RX time: $t_{\text{RX}} = \SI{50}{\milli\second}$ at $I_{\text{RX}} = \SI{3.9}{\milli\ampere}$.
%   \item Total active time per wake-up:
%   \begin{equation}
%     t_{\text{act}} = t_{\text{WU}} + t_{\text{ADC}} + t_{\text{TX}} + t_{\text{RX}}
%     = \SI{161.83}{\milli\second}.
%   \end{equation}
%   \item EM0 current consumption at $f=\SI{76.8}{\mega\hertz}$ with $\SI{27}{\micro\ampere\per\mega\hertz}$:
%   \begin{equation}
%     I_{\text{EM0}} = \SI{27}{\micro\ampere\per\mega\hertz}\cdot \SI{76.8}{\mega\hertz}
%     = \SI{2.0736}{\milli\ampere}.
%   \end{equation}
%   \item EM4 sleep current: $I_{\text{EM4}} = \SI{0.17}{\micro\ampere}$.
% \end{itemize}


% \textbf{Charge consumption per wake-up}\newline
% The active-mode charge drawn during wake-up and ADC sampling is approximated as:
% \begin{equation}
%   Q_{\text{EM0}} = (t_{\text{WU}}+t_{\text{ADC}})\cdot I_{\text{EM0}}
%   = (\SI{1.83}{\milli\second}+\SI{10}{\milli\second}) \cdot \SI{2.0736}{\milli\ampere}
%   = \SI{0.0245306888}{\milli\ampere\second}.
% \end{equation}

% The charge drawn during Tx and Rx is:
% \begin{equation}
%   Q_{\text{TX}} = t_{\text{TX}}\cdot I_{\text{TX}}
%   = \SI{100}{\milli\second}\cdot \SI{8.2}{\milli\ampere}
%   = \SI{0.82}{\milli\ampere\second},
% \end{equation}
% \begin{equation}
%   Q_{\text{RX}} = t_{\text{RX}}\cdot I_{\text{RX}}
%   = \SI{50}{\milli\second}\cdot \SI{3.9}{\milli\ampere}
%   = \SI{0.195}{\milli\ampere\second}.
% \end{equation}

% Thus, the total charge per wake-up is:
% \begin{equation}
%   Q_{\text{WU}} = Q_{\text{EM0}}+Q_{\text{TX}}+Q_{\text{RX}}
%   = \SI{1.0395306888}{\milli\ampere\second}.
% \end{equation}

% Converted to \si{\milli\ampere\hour}:
% \begin{equation}
%   C_{\text{WU}} = \frac{Q_{\text{WU}}}{\SI{3600}{\second\per\hour}}
%   = \frac{\SI{1.0395306888}{\milli\ampere\second}}{\SI{3600}{\second\per\hour}}
%   = \SI{2.8876e-4}{\milli\ampere\hour}.
% \end{equation}

% \textbf{Active consumption over 10 years}\newline
% The total active consumption over $N_{\text{WU}}$ wake-ups becomes:
% \begin{equation}
%   C_{\text{act,10y}} = N_{\text{WU}}\cdot C_{\text{WU}}
%   = \SI{123}{}\cdot \SI{2.8876e-4}{\milli\ampere\hour}
%   = \SI{0.03552}{\milli\ampere\hour}.
% \end{equation}

% \textbf{Sleep consumption (EM4) over 10 years}\newline
% The sleep consumption is approximated as:
% \begin{equation}
%   C_{\text{sleep,10y}} = I_{\text{EM4}}\cdot H_{10y}
%   = \SI{0.17}{\micro\ampere}\cdot \SI{87600}{\hour}
%   = \SI{14.892}{\milli\ampere\hour}.
% \end{equation}

% \textbf{Total consumption}\newline
% The total estimated consumption is:
% \begin{equation}
%   C_{\text{tot,10y}} = C_{\text{act,10y}} + C_{\text{sleep,10y}}
%   = \SI{0.03552}{\milli\ampere\hour} + \SI{14.892}{\milli\ampere\hour}
%   \approx \SI{14.9275}{\milli\ampere\hour}.
% \end{equation}
% Notice this i not taking other risks into account such as current leakage, rejoin delays and temperature impacts.

 

